Put \(c=1/e\).  Affine additions \(b(t-1)\) leave both the divergence and the calibrated radius invariant because every likelihood ratio has \(Q\)-mean one and \(e(c-1)-(1-e)=0\).  Convexity gives the chord bound
\[
 f(t)\le(1-t/c)f(0)+(t/c)f(c),\qquad 0\le t\le c. \tag{1}
\]
Every \(Q=eP+(1-e)R\) with \(R\ll P\) satisfies \(Q\sim P\) and \(dP/dQ\le c\).  Integration of (1) therefore gives \(D_f(P\|Q)\le B_f(e)\), proving \(\mathfrak L_e(P)\subseteq\mathfrak B_{f,e}(P)\).

If \(f\in\mathfrak H_c\), choose an affine representative \(g\) that vanishes on \([0,c]\) and is strictly positive above \(c\).  Then \(B_g(e)=0\).  For \(Q\in\mathfrak B_{f,e}(P)\), with \(r=dP/dQ\),
\[
 0\le\int g(r)\,dQ\le0,
\]
so \(r\le c\), \(Q\)-almost surely.  The probability law \(R=(Q-eP)/(1-e)\) is nonnegative by this cap and satisfies \(R\ll P\) because \(Q\ll P\).  Hence \(Q\in\mathfrak L_e(P)\), proving universal exactness.

For necessity, assume exactness for every full-support binary \(P=\operatorname{Bernoulli}(p)\).  The boundary \(Q_{ep}=\operatorname{Bernoulli}(ep)\) lies in \(\mathfrak L_e(P)\), since its residual \(\operatorname{Bernoulli}(0)\) is absolutely continuous with respect to full-support \(P\); every \(Q_q=\operatorname{Bernoulli}(q)\) with \(0<q<ep\) is outside because \(p/q>c\).  Continuity from the exterior gives
\[
 epf(c)+(1-ep)f\!\left(\frac{1-p}{1-ep}\right)=B_f(e). \tag{2}
\]
Letting \(p\downarrow0\) yields \(B_f(e)=0\).  Normalize \(g(t)=f(t)+f(0)(t-1)\); then \(g(0)=g(1)=g(c)=0\), and (2) gives \(g=0\) on \([0,1]\).  For \(y\in(1,c)\), combine the ratios \(y\) and any \(x\in(0,1)\) with weight \(\lambda=(1-x)/(y-x)\).  The resulting full-support binary pair has mean likelihood ratio one and respects the cap, so exactness and convexity at the mean-one likelihood ratio force
\[
 0\le\lambda g(y)+(1-\lambda)g(x)\le B_g(e)=0.
\]
Thus \(g=0\) on \([0,c]\).  Convexity gives \(g(t)\ge0\) above \(c\); equality at any \(t>c\) would make the analogous two-ratio full-support binary pair a zero-divergence ball member violating the cap.  Hence \(g(t)>0\) for all \(t>c\), so \(f\in\mathfrak H_c\).  The bow-mixture completeness theorem identifies the exact law class with the compatible causal class, completing the three-way equivalence.